\chapter{Special Relativity}

\section{Relativistic Momentum}

\section{Relativistic Energy}

The total energy \(E\) of an object moving at relativistic speed \(v\) is
\begin{equation} \label{eq:total-relativistic-energy}
    E = \frac{m c^2}{\sqrt{1 - \frac{v^2}{c^2}}} = \gamma m c^2,
\end{equation}
where
\(m\) is the mass of the object,
\(c \coloneq \qty{299792458}{\metre \per \second}\) is the speed of
light in a vacuum,
and
\begin{equation}
    \gamma \coloneq \frac{1}{\sqrt{1 - \frac{v^2}{c^2}}}.
\end{equation}

From \pgRef{eq:total-relativistic-energy},
we can conclude that the energy \(E_\mathrm{S}\) of a stationary object
(i.e., \(v = \qty{0}{\metre \per \second}\))
is
\begin{equation}
    E_\mathrm{S} = \frac{m c^2}{\sqrt{1 - \frac{0}{c^2}}} = m c^2.
\end{equation}
Thus, the kinetic energy \(E_\mathrm{K}\) is
\begin{equation}
    E_\mathrm{K}
    = E - E_\mathrm{S}
    = m c^2 (\gamma - 1).
\end{equation}

\begin{subequations}
    In the limit where \(v \ll c\),
    we can use the binomial approximation to find
    \begin{equation}
        \begin{aligned}
            \gamma
            = \frac{1}{\sqrt{1 - \frac{v^2}{c^2}}}
            &= \ab(1 - \frac{v^2}{c^2})^{-\frac{1}{2}}
            &\approx 1 + \frac{v^2}{2c^2}.
        \end{aligned}
    \end{equation}
    Therefore, in the classical limit,
    \begin{equation}
        E_\mathrm{K}
        \approx m c^2 (1 + \frac{v^2}{2c^2} - 1)
        = m v^2
    \end{equation}
\end{subequations}

Furthermore,the relationship between total energy \(E\) and momentum \(p\) is
\begin{equation}
    E = p c^2 + (m c^2)^2.
\end{equation}
